\section{Quadratic path costs}\label{qp:section}

\subsection{The quotient norm and its normalization}

For an integer $n\ge1$, let $T_n=\mathbb N^{\le n}$, including the empty
sequence $o$, with the initial-segment order $\preceq$.  Work over the
real field.  In $H_n=\ell_2(T_n)$ put
\[
 p_t=\sum_{s\preceq t}e_s,\qquad
 J_n\lambda=\sum_{t\in T_n}\lambda_t p_t
       \quad(\lambda\in\ell_1(T_n)).
\]
The series converges absolutely in $H_n$, since
$\|p_t\|_2\le\sqrt{n+1}$.  Thus
$(J_n\lambda)(s)=\sum_{t\succeq s}\lambda_t$ and
$\|J_n\|\le\sqrt{n+1}$.  Define $Z_n$ to be $H_n$ with the norm
\begin{equation}\label{qp:norm}
 |z|_n=\inf_{z=h+J_n\lambda}
       \bigl(\|h\|_2^2+\|\lambda\|_1^2\bigr)^{1/2},
 \qquad h\in H_n,\quad\lambda\in\ell_1(T_n),
\end{equation}
and set
\[
 X_{\mathrm q}=\Bigl(\bigoplus_{n\ge1}Z_n\Bigr)_2,
 \qquad \|z\|_{\mathrm q}^2=\sum_n|z_n|_n^2.
\]
In the notation of Section~\ref{sec:linear-paths}, $p_t=b_t$,
$J_n=B_n$, and $|z|_n=Q_n(z)$.
To bound midpoint tails, we compare the costs of two endpoint
representations with the cost of their average restricted to a finite head.

\begin{proposition}\label{qp:construction}
The spaces $Z_n$ and $X_{\mathrm q}$ are Banach and reflexive, and finite
node support is dense in $X_{\mathrm q}$.  For $z\in H_n$,
\begin{equation}\label{qp:comparison}
 \frac{\|z\|_2}{\sqrt{n+2}}\le |z|_n\le\|z\|_2,
 \qquad |z_t|\le\sqrt2\,|z|_n.
\end{equation}
In particular $1/\sqrt2\le |p_t|_n,|e_t|_n\le1$, and
$|e_o|_n=1/\sqrt2$.  For $f\in\ell_2(T_n)$ the dual norm is
\begin{equation}\label{qp:dual}
 |f|_{n,*}=
 \left(\|f\|_2^2+
     \sup_{t\in T_n}\left|\sum_{s\preceq t}f_s\right|^2\right)^{1/2}.
\end{equation}
\end{proposition}

\begin{proof}
Scaling and adding representations in \eqref{qp:norm} proves
homogeneity and the triangle inequality.  Every representation satisfies
\[
 \|z\|_2\le\|h\|_2+\sqrt{n+1}\|\lambda\|_1
 \le\sqrt{n+2}\bigl(\|h\|_2^2+\|\lambda\|_1^2\bigr)^{1/2};
\]
the choice $h=z$ gives the upper bound.  This proves positive
definiteness, completeness, and equivalence to the Hilbert norm.
The dual and bidual therefore have the same underlying continuous
linear functionals as for $H_n$, so $Z_n$ is reflexive.  For each
coordinate,
$|z_t|\le\|h\|_2+\|\lambda\|_1$, which gives the second assertion of
\eqref{qp:comparison}.  A single atomic mass represents $p_t$ with cost
one.  For $e_t$ use its Hilbert representation.  At the root, the
representation $e_o=e_o/2+J_n(\delta_o/2)$ attains the matching cost
$1/\sqrt2$.

The quotient map $(h,\lambda)\mapsto h+J_n\lambda$ has domain
$H_n\oplus_2\ell_1(T_n)$.  Pulling the functional $f$ back gives
$(f,J_n^*f)$, whose norm is
$\bigl(\|f\|_2^2+\|J_n^*f\|_\infty^2\bigr)^{1/2}$.
Indeed the upper bound is Cauchy--Schwarz, and the reverse bound follows
by choosing the Hilbert vector in the direction of $f$, an atomic
coordinate arbitrarily close to the supremum for $J_n^*f$, and then
optimizing their two nonnegative lengths.  The pullback and quotient
functional norms agree: a representation bounds the functional by its
cost, and representations can have cost arbitrarily close to
$|z|_n$.  Since $(J_n^*f)(t)=\sum_{s\preceq t}f_s$, this proves
\eqref{qp:dual} without requiring a maximizing endpoint.

For completeness of the outer sum, a Cauchy sequence converges in each
$Z_n$; its Cauchy bounds pass first to every finite sum and then to the
whole sum, proving membership and convergence of the limit.  A bounded
functional on this sum restricts to functionals $f_n\in Z_n^*$ with
$\sum_n\|f_n\|^2<\infty$.  To prove the bound on this sum, test on
finite collections of almost norming unit vectors multiplied by
arbitrary Euclidean unit scalar vectors.  Conversely such $(f_n)$
defines a functional by Cauchy--Schwarz.  Density of finite component
sums shows that these constructions are inverse and isometric.
Applying the same description to the bidual and using component
reflexivity proves reflexivity of $X_{\mathrm q}$.  Finally truncate
the outer sum and then the Hilbert coordinates in each remaining
component to obtain finite node support.
\end{proof}

Write
\[
 H=\Bigl(\bigoplus_n H_n\Bigr)_2,\qquad
 L=\Bigl(\bigoplus_n\ell_1(T_n)\Bigr)_2.
\]
The componentwise map $J:L\to X_{\mathrm q}$ and the coordinate
inclusion $H\to X_{\mathrm q}$ are contractions, and
\begin{equation}\label{qp:global-quotient}
 \|z\|_{\mathrm q}=
 \inf_{z=h+J\lambda}
       \bigl(\|h\|_H^2+\|\lambda\|_L^2\bigr)^{1/2}.
\end{equation}
For the reverse inequality in this formula choose component
representations with total error in the squared costs less than any
prescribed positive number.  This also produces $h\in H$ and
$\lambda\in L$.  The map $J$ here takes values in $X_{\mathrm q}$; no
uniform bound $L\to H$ is used.

A \emph{finite ancestral head} is a finite set $A$ of nodes in the
disjoint union of the $T_n$, closed under taking ancestors within each
component. This is the closure condition called initial in the preceding
sections. Its intersection $A_n$ with $T_n$ may be empty. Let
$P_A$ retain the coordinates in $A$ and let $Q_A=I-P_A$.  For
$t\in T_n$ denote its longest ancestor in $A_n$ by $r(t)$, using
$\perp$ when there is none, and put $p_\perp=0$.

\begin{lemma}\label{qp:head}
The operator $R_A$ which sends the mass $\lambda_t$ to $r(t)$ and
discards it when $r(t)=\perp$ is contractive on $L$.  Moreover
\[
 P_AJ=JR_A,\qquad \|P_A\|\le1,\qquad \|Q_A\|\le2,
 \qquad |Q_AJ_n\lambda|_n\le2\|\lambda\|_1.
\]
Here the last formula is in a single component.  In particular the
path $p_t-p_{r(t)}$ outside the head has norm at most two.
\end{lemma}

\begin{proof}
The triangle inequality bounds the absolute sum in every fiber of
$r$, so $R_A$ contracts each $\ell_1$ component and hence $L$.
Coordinate restriction gives $P_Ap_t=p_{r(t)}$ and thus $P_AJ=JR_A$.
Projecting $h$ orthogonally in $H$ and folding $\lambda$ in
\eqref{qp:global-quotient} proves contractivity of $P_A$.
The remaining bounds follow from $Q_A=I-P_A$ and
$\|J\lambda\|_{\mathrm q}\le\|\lambda\|_L$.
\end{proof}

\begin{proposition}\label{qp:no-auc}
No equivalent norm on $X_{\mathrm q}$ is asymptotically uniformly
convex.  More precisely, if
$a\|z\|_{\mathrm q}\le N(z)\le b\|z\|_{\mathrm q}$, with
$0<a\le b<\infty$, then
$\overline\delta_N(t_0)=0$ at $t_0=a/(\sqrt2 b)$.
\end{proposition}

\begin{proof}
For each fixed nonterminal $s\in T_n$, the coordinate vectors
$e_{s^\frown j}$ are weakly null in $Z_n$, by its equivalence to
$H_n$, and therefore in $X_{\mathrm q}$.  Their norms and the norms
of $p_s$ are at least $1/\sqrt2$, while $\|p_s\|_{\mathrm q}\le1$.
Suppose $\overline\delta_N(t_0)>0$ and choose
$c>0$ with $2c<\overline\delta_N(t_0)$.  At
$u=p_s/N(p_s)$ choose a closed finite-codimensional subspace $F$
such that $N(u+t_0v)>1+2c$ for every $v\in F$ with $N(v)=1$.
The normalized sequence $v_j=e_{s^\frown j}/N(e_{s^\frown j})$
is weakly null, because its denominators are bounded below.  Its
image in $X_{\mathrm q}/F$ is norm null.  Approximation in $F$,
followed by normalization, gives vectors of $F\cap S_N$ whose
$N$-distance from $v_j$ tends to zero.  Consequently
$N(u+t_0v_j)>1+c$ for all sufficiently large $j$.

For a convex function $r\mapsto N(u+rv_j)$ with value one at zero,
a lower bound greater than one at $t_0$ persists for $r\ge t_0$.
Apply this with
$r=N(e_{s^\frown j})/N(p_s)\ge a/(\sqrt2 b)$ to obtain
$N(p_{s^\frown j})\ge(1+c)N(p_s)$.
Choosing successors for $n$ steps gives
$b\ge (1+c)^n a/\sqrt2$, impossible for arbitrarily large $n$.
The asymptotic modulus is nonnegative: for each unit center one may
restrict to the kernel of a norming functional.  It is therefore zero.
\end{proof}

\subsection{What is lost by averaging and restricting}

Fix a finite ancestral head $A$, write $P=P_A$, $Q=Q_A$, and suppose
$Px=x$.  Choose representations
\[
 x\pm y=h^\pm+J\lambda^\pm.
\]
We will always choose finite total squared cost, as allowed by
\eqref{qp:global-quotient}.  Put
\[
 \bar h=\tfrac12(h^++h^-),\quad
 \widetilde h=\tfrac12(h^+-h^-),\quad
 \sigma=\tfrac12(\lambda^++\lambda^-),\quad
 \tau=\tfrac12(\lambda^+-\lambda^-),
\]
and, in each component, put
\[
 M_n=\tfrac12(\|\lambda_n^+\|_1+\|\lambda_n^-\|_1),\qquad
 \nu_n=(R_A\sigma)_n,\qquad q_n=\|\nu_n\|_1,\qquad
 d_n=M_n-q_n.
\]
Thus $x=P\bar h+J\nu$ and $d_n\ge0$. To separate the two sources of
this loss, view
$m_n=(|\lambda_n^+|+|\lambda_n^-|)/2$ as a nonnegative endpoint
measure and $\sigma_n$ as a signed endpoint measure. Let
$C_r=\{t:r(t)=r\}$ for $r\in A_n\cup\{\perp\}$. Absolute convergence
gives
\[
 d_n=\sum_{r\in A_n}\bigl(m_n(C_r)-|\sigma_n(C_r)|\bigr)
          +m_n(C_\perp).
\]
Each term is nonnegative. The first sum records cancellation during
averaging and folding within retained fibres; the last term records
the mass discarded when the head is empty in that component.

\begin{lemma}\label{qp:energy}
Let
\[
 C=\frac12\sum_{\epsilon\in\{+,-\}}
       (\|h^\epsilon\|_H^2+\|\lambda^\epsilon\|_L^2),
 \qquad E=C-\|x\|_{\mathrm q}^2.
\]
Then
\begin{align}
 \|Q\bar h\|_H^2+\|\widetilde h\|_H^2+
       \sum_n(M_n^2-q_n^2)&\le E,\label{qp:energy-mean}\\
 \tfrac12(\|Qh^+\|_H^2+\|Qh^-\|_H^2)+
       \sum_n d_n^2&\le E.\label{qp:energy-endpoints}
\end{align}
\end{lemma}

\begin{proof}
The Hilbert parallelogram identity and the scalar square inequality
give
$C\ge\|\bar h\|_H^2+\|\widetilde h\|_H^2+\sum_nM_n^2$.
On the other hand the displayed representation of $x$ has squared
cost $\|P\bar h\|_H^2+\sum_nq_n^2\ge\|x\|_{\mathrm q}^2$.
Subtract, using the orthogonal decomposition in $H$, to prove
\eqref{qp:energy-mean}.  Now
$M_n^2-q_n^2\ge(M_n-q_n)^2$ and
\[
 \tfrac12(\|Qh^+\|_H^2+\|Qh^-\|_H^2)
 =\|Q\bar h\|_H^2+\|Q\widetilde h\|_H^2.
\]
These identities imply \eqref{qp:energy-endpoints}.
\end{proof}

To convert this loss into a bound for the vector tail, we need an
elementary representation of a positive array on the outside forest.
Its importance is that the bound is independent of the height and
allows countably many children.

\begin{lemma}[Stopped flows]\label{qp:flow}
Let $O=T_n\setminus A_n$, let $\mathcal R$ be the roots of its
connected subtrees, and suppose $V:O\to[0,\infty)$ satisfies
\[
 V(t)\ge\sum_{c\text{ child of }t}V(c),\qquad
 B:=\sum_{c\in\mathcal R}V(c)<\infty.
\]
For any signs $\varepsilon_t$ constant on each connected subtree of
$O$, the vector $\varepsilon_tV(t)$ on $O$, extended by zero on
$A_n$, has norm at most $2B$ in $Z_n$.  The pointwise minimum of two
arrays satisfying the child inequality also satisfies it.
\end{lemma}

\begin{proof}
Set $\alpha_t=V(t)-\sum_{c\text{ child of }t}V(c)\ge0$.
Induction from the last level, using nonnegative sums, gives
$V(t)=\sum_{s\succeq t}\alpha_s$ and
$\sum_{t\in O}\alpha_t=B$.  This induction is over the finite
height, not over a finite set of children.  Hence the required vector is
\[
 \sum_{t\in O}\varepsilon_t\alpha_t
                      (p_t-p_{r(t)}).
\]
The series is absolutely convergent in $Z_n$ by Lemma~\ref{qp:head}
and has norm at most $2B$.  Continuity of the coordinates identifies
its sum with the stated array.  Finally
\[
 \sum_c\min(V(c),W(c))
 \le\min\Bigl(\sum_cV(c),\sum_cW(c)\Bigr)
 \le\min(V(t),W(t)),
\]
which proves the assertion about minima.
\end{proof}

\subsection{Clipping one endpoint against the minority mass}

\begin{theorem}\label{thm:quadratic-tail}
Let $A$ be a finite ancestral head in $X_{\mathrm q}$ and let
$P_Ax=x$.  For every $y\in X_{\mathrm q}$ and $R\ge0$,
\begin{equation}\label{qp:tail-eight}
 \|x+y\|_{\mathrm q},\ \|x-y\|_{\mathrm q}\le R
 \quad\Longrightarrow\quad
 \|Q_Ay\|_{\mathrm q}\le8\sqrt{R^2-\|x\|_{\mathrm q}^2}.
\end{equation}
The displacement $y$ need not vanish on the head.
\end{theorem}

\begin{proof}
Convexity gives $\|x\|_{\mathrm q}\le R$, so the right side is
defined.  The case $R=0$ is immediate.  Choose the representations
above with each squared cost at most $R^2+\eta$, where $\eta>0$.
Then $E\le D:=R^2+\eta-\|x\|_{\mathrm q}^2$.

Fix a component.  For every fiber of $r(t)$ choose its majority sign
$\varepsilon_r$ among the coefficients of both labeled arrays
$\lambda^+,\lambda^-$.  If the positive and negative masses in a
nondiscarded fiber are $a_r,b_r$, its contribution to $d_n$ is
$\min(a_r,b_r)$; for a discarded fiber its contribution is
$(a_r+b_r)/2\ge\min(a_r,b_r)$.  Thus the total minority mass over
the two arrays is at most $d_n$.

Outside $A_n$, each descendant has the same $r(t)$ as its ancestor.
For $\epsilon\in\{+,-\}$ define the majority and minority
occupancies
\[
 U^\epsilon(t)=\sum_{s\succeq t}(\varepsilon_{r(s)}
                                    \lambda_n^\epsilon(s))_+,
 \qquad
 V^\epsilon(t)=\sum_{s\succeq t}(-\varepsilon_{r(s)}
                                    \lambda_n^\epsilon(s))_+,
 \qquad V=V^++V^-.
\]
Here $a_+=\max(a,0)$.  All these arrays satisfy the child inequality
of Lemma~\ref{qp:flow}, and the total initial mass of $V$ is at most
$d_n$.  Since $x$ vanishes outside the head,
\[
 \varepsilon_{r(t)}(U^+(t)+U^-(t)-V(t))
                       =-h_n^+(t)-h_n^-(t).
\]
Let $S(t)=\min(U^+(t),V(t))$.  It is a stopped flow with initial
mass at most $d_n$, and
\[
 0\le U^+(t)-S(t)
 \le (U^+(t)+U^-(t)-V(t))_+
 \le |h_n^+(t)+h_n^-(t)|.
\]
Consequently the tail of $y_n$, which agrees with that of $x_n+y_n$,
is the sum of the signed flows $S,-V^+$ and the Hilbert vector
$h_n^++\varepsilon_r(U^+-S)$, all restricted to the outside.
Lemma~\ref{qp:flow} gives
\[
 |Q_Ay_n|_n\le4d_n+\|Q_Ah_n^+\|_2+
                              \|Q_A(h_n^++h_n^-)\|_2.
\]
Take the outer Euclidean norm and use
\eqref{qp:energy-endpoints}.  The three terms are at most
$4\sqrt D$, $\sqrt{2D}$, and $2\sqrt D$, respectively.
Thus $\|Q_Ay\|_{\mathrm q}\le(6+\sqrt2)\sqrt D\le8\sqrt D$.
Let $\eta\downarrow0$.  No minimizing representation was assumed.
\end{proof}

The next corollary quantifies the symmetric-lens geometry in the
characterization of AMUC by Dilworth et al.~\cite[Theorem~2.1]{DKRRZ2016}.

\begin{corollary}\label{qp:lens-eight}
If $x\in X_{\mathrm q}$ and an infinite sequence $(y_j)$ satisfies
$\|x\pm y_j\|_{\mathrm q}\le R$ and
$\|y_i-y_j\|_{\mathrm q}\ge\varepsilon>0$ for $i\ne j$, then
\[
 \|x\|_{\mathrm q}^2+\frac{\varepsilon^2}{256}\le R^2.
\]
\end{corollary}

\begin{proof}
Approximate $x$ within $\rho>0$ by $x_0$ supported in a finite
ancestral head $A$.  Theorem~\ref{thm:quadratic-tail} bounds each
$Q_Ay_j$ by $8\sqrt{(R+\rho)^2-\|x_0\|_{\mathrm q}^2}$.
The sequence $(y_j)$ is bounded, and $P_A$ has finite rank, so two
distinct head projections can be arbitrarily close.  Their difference
and the two tail bounds give
$\varepsilon\le16\sqrt{(R+\rho)^2-\|x_0\|_{\mathrm q}^2}$.
Let $\rho\downarrow0$.
\end{proof}

The arbitrary-tail estimate controls the complement of a head without
requiring the displacement to vanish on the head. The remaining arguments
organize the same cancellation budget for other purposes: symmetric
clipping keeps the average of the two squared endpoint norms; stopping
and majority pruning retain different parts of the signed mass; and a
componentwise passage produces a finite-codimensional midpoint estimate.

\subsection{A symmetric clipping estimate}

The proof clips both endpoint arrays while keeping the two Hilbert
variables separate.

\begin{proposition}\label{qp:symmetric-twelve}
For a finite ancestral head $A$, $P_Ax=x$, and $P_Ay=0$,
\begin{equation}\label{qp:symmetric-144}
 \|x\|_{\mathrm q}^2+\frac{\|y\|_{\mathrm q}^2}{144}
 \le\frac{\|x+y\|_{\mathrm q}^2+\|x-y\|_{\mathrm q}^2}{2}.
\end{equation}
\end{proposition}

\begin{proof}
It suffices to prove the assertion in one component, since the
result can then be summed.  Choose endpoint representations whose
squared costs exceed $|x\pm y|_n^2$ by at most $\eta>0$, and use
the notation of Lemma~\ref{qp:energy}, omitting the component index.
Put $E=(|x+y|_n^2+|x-y|_n^2)/2+\eta-|x|_n^2$.
Then
\[
 \|Q\bar h\|_2^2+\|\widetilde h\|_2^2+M^2-q^2\le E.
\]
Orient each nondiscarded fiber by the sign of its folded mean
coefficient $\nu_{r(t)}$, choosing $+1$ when it is zero, and choose
either sign in $\{-1,1\}$ on a discarded fiber.
Define half-weighted positive and negative endpoint masses by
\[
 \alpha_t^\epsilon=\tfrac12(\varepsilon_{r(t)}
                         \lambda_t^\epsilon)_+,
 \qquad
 \beta_t^\epsilon=\tfrac12(-\varepsilon_{r(t)}
                         \lambda_t^\epsilon)_+,
 \quad\epsilon\in\{+,-\}.
\]
Let $\beta=\beta^++\beta^-$.  In a nondiscarded fiber the
positive mass minus negative mass is the absolute folded
coefficient.  The discarded fiber contributes zero to $q$.
It follows that $\|\beta\|_1\le M-q\le\sqrt E$.

For $t\notin A_n$ put
$A^\epsilon(t)=\sum_{s\succeq t}\alpha_s^\epsilon$ and
$B^\epsilon(t)=\sum_{s\succeq t}\beta_s^\epsilon$,
and let $A_{\mathrm{tot}}=A^++A^-$,
$B_{\mathrm{tot}}=B^++B^-$. The endpoint identities give
\[
 0=\bar h(t)+\varepsilon_{r(t)}(A_{\mathrm{tot}}(t)-B_{\mathrm{tot}}(t)),
 \qquad
 y(t)=\widetilde h(t)+\varepsilon_{r(t)}
          (A^+(t)-A^-(t)-B^+(t)+B^-(t)).
\]
For each sign put $S^\epsilon=\min(A^\epsilon,B_{\mathrm{tot}})$ and
$V^\epsilon=A^\epsilon-S^\epsilon$.  Then
$0\le V^\epsilon\le(A_{\mathrm{tot}}-B_{\mathrm{tot}})_+
\le|\bar h|$ outside the head.
Each $S^\epsilon$ is a stopped flow of initial mass at most
$\|\beta\|_1$, and the same bound applies to each
$B^\epsilon$.  Lemma~\ref{qp:flow} estimates their four signed
vectors by $8\|\beta\|_1$ in total.  The remaining terms have
Hilbert norm at most $\|Q\widetilde h\|_2+2\|Q\bar h\|_2$.
Since $y$ vanishes on the head,
\[
 |y|_n\le\|Q\widetilde h\|_2+2\|Q\bar h\|_2
                 +8\|\beta\|_1\le12\sqrt E.
\]
Let $\eta\downarrow0$, and sum the resulting component
inequalities.  Empty heads cause no exception: their entire mass
is in the discarded fiber, and the same bounds hold.
\end{proof}

\begin{corollary}\label{qp:lens-twelve}
If $\|x\pm y_j\|_{\mathrm q}\le r$ and
$\|y_i-y_j\|_{\mathrm q}\ge2d>0$ for $i\ne j$ in an infinite
sequence, then
\[
 \|x\|_{\mathrm q}^2+d^2/144\le r^2.
\]
\end{corollary}

\begin{proof}
Choose $x_0$ of finite node support with
$\|x-x_0\|_{\mathrm q}<\epsilon$, and a finite ancestral head
$A$ supporting it.  Finite rank and boundedness give distinct
$i,j$ with $\|P_Az\|_{\mathrm q}<\epsilon$ for
$z=(y_i-y_j)/2$.  Convexity gives $\|x\pm z\|_{\mathrm q}\le r$,
while $\|z\|_{\mathrm q}\ge d$.  Apply
\eqref{qp:symmetric-144} to $x_0,Q_Az$.  These vectors have endpoint
norms at most $r+2\epsilon$, and
$\|Q_Az\|_{\mathrm q}\ge d-\epsilon$.  Let $\epsilon\downarrow0$.
\end{proof}

\subsection{Stopping where signed mass cancels}

Instead of orienting the coefficients, one can stop where their
signed sum is less than half of their absolute mass.  The discarded
subtrees are disjoint, so their total mass is paid for by the same
loss $d_n$.  This gives a smaller numerical constant by a different
decomposition.

\begin{proposition}\label{qp:tail-seven}
Under the hypotheses of Theorem~\ref{thm:quadratic-tail},
\[
 \|Q_Ay\|_{\mathrm q}\le7\sqrt{R^2-\|x\|_{\mathrm q}^2}.
\]
\end{proposition}

\begin{proof}
Choose the endpoint representations with squared costs at most
$(R+\eta)^2$.  In the notation preceding Lemma~\ref{qp:energy}
define the nonnegative endpoint measure
$m_n=(|\lambda_n^+|+|\lambda_n^-|)/2$.
Both $|\sigma_n|$ and $|\tau_n|$ are bounded pointwise by $m_n$,
and $\|m_n\|_1=M_n$.  The lemma gives
\begin{equation}\label{qp:seven-energy}
 \|Q\bar h\|_H^2+\|\widetilde h\|_H^2+\sum_nd_n^2
 \le (R+\eta)^2-\|x\|_{\mathrm q}^2=:E.
\end{equation}

Write $S_t=\{s:s\succeq t\}$ for a descendant set.  Mark a node
$t\notin A_n$ when
$|\sigma_n(S_t)|<m_n(S_t)/2$, and retain only the first marked
nodes on each path outside the head.  Every marked node has such
a first ancestor, since it has finitely many ancestors.  The
retained descendant sets are pairwise disjoint.

Each retained set lies in a single fiber of $r$.  If disjoint
subsets $S_j$ lie in a nondiscarded fiber $C$, then
\[
 \sum_j\bigl(m_n(S_j)-|\sigma_n(S_j)|\bigr)
 \le m_n(C)-|\sigma_n(C)|.
\]
Indeed $|\sigma_n(C)|$ is at most the sum of the displayed absolute
sums plus the $m_n$-mass of the remaining set.  This argument also
holds for countably many subsets by absolute convergence.
Discarded fibers are bounded by their whole $m_n$-mass.  Summing
over fibers yields
\[
 \sum_{t\text{ first marked}}
      \bigl(m_n(S_t)-|\sigma_n(S_t)|\bigr)\le d_n,
 \qquad
 \sum_{t\text{ first marked}}m_n(S_t)\le2d_n.
\]
Split $\tau_n=\tau_n^{\mathrm r}+\tau_n^{\mathrm k}$ by
restricting the first part to the union of these descendant sets.
Then $\|\tau_n^{\mathrm r}\|_1\le2d_n$ and
\[
 \|QJ\tau^{\mathrm r}\|_{\mathrm q}
          \le4\Bigl(\sum_nd_n^2\Bigr)^{1/2}.
\]
If an outside node has a first marked ancestor, the occupancy of
$\tau_n^{\mathrm k}$ there is zero.  Otherwise it is unmarked, and
\[
 |(J_n\tau_n^{\mathrm k})(t)|
 \le m_n(S_t)\le2|\sigma_n(S_t)|=2|\bar h_n(t)|.
\]
The last equality uses $Qx=0$ in $x=\bar h+J\sigma$.
Thus $QJ\tau^{\mathrm k}\in H$ and has Hilbert norm at most
$2\|Q\bar h\|_H$.  Since $y=\widetilde h+J\tau$,
\[
 \|Qy\|_{\mathrm q}
 \le\|Q\widetilde h\|_H+2\|Q\bar h\|_H
           +4\Bigl(\sum_nd_n^2\Bigr)^{1/2}
 \le7\sqrt E.
\]
Let $\eta\downarrow0$.
\end{proof}

\begin{corollary}\label{qp:lens-seven}
For $0<c\le1$ put $\delta=c^2/2000$.  If
$\|x\pm y_j\|_{\mathrm q}\le1$ and
$\|y_i-y_j\|_{\mathrm q}\ge2c$ for all distinct indices of an
infinite sequence, then $\|x\|_{\mathrm q}\le1-\delta$.
\end{corollary}

\begin{proof}
If $\|x\|_{\mathrm q}>1-\delta$, choose a finite-support $x_0$
within $\delta$ of $x$, and an ancestral head $A$ for $x_0$.
Then $\|x_0\|_{\mathrm q}>1-2\delta$ and
$\|x_0\pm y_j\|_{\mathrm q}\le1+\delta$.
Proposition~\ref{qp:tail-seven} gives
\[
 \|Q_Ay_j\|_{\mathrm q}
 <7\sqrt{(1+\delta)^2-(1-2\delta)^2}
 \le7\sqrt{6\delta}<c/2.
\]
Two bounded head projections are less than $c$ apart, contradicting
the assumed separation of the full vectors.
\end{proof}

\subsection{Pruning by a majority sign}

This argument tracks the labeled majority terms before cancellation in
the half-difference. It keeps those terms until their occupancy is
comparable with the minority occupancy, separating removed mass from
a remaining occupancy bounded in Hilbert norm. This decomposition will
also be used after approximating an arbitrary center.

\begin{proposition}\label{qp:tail-nine}
Under the hypotheses of Theorem~\ref{thm:quadratic-tail},
\[
 \|Q_Ay\|_{\mathrm q}\le9\sqrt{R^2-\|x\|_{\mathrm q}^2}.
\]
\end{proposition}

\begin{proof}
Use representations with squared costs at most $R^2+\eta$ and
write $E=R^2+\eta-\|x\|_{\mathrm q}^2$.
In each nondiscarded fiber $C_r$, orient the coefficients by the sign of
the folded coefficient $\nu_n(r)$, choosing $+1$ when it is zero; use either
sign in $\{-1,1\}$ in a discarded fiber.  Let $g_n,b_n$ be
half the total absolute masses of coefficients agreeing and
disagreeing with this sign, respectively, counting both labeled
endpoint representations.  Then $\|b_n\|_1\le d_n$.  More
precisely, if $A_n\ne\varnothing$, all nodes have an ancestor in
$A_n$ and $2\|b_n\|_1=d_n$; when $A_n=\varnothing$, simply
$d_n=M_n$.

On the outside put $G(t)=g_n(S_t)$ and $B(t)=b_n(S_t)$.  Constancy
of the orientation along each outside subtree gives
\[
 |(J_n\sigma_n)(t)|=|G(t)-B(t)|=|\bar h_n(t)|.
\]
In the labeled half-difference
$\tau_n=(\lambda_n^+-\lambda_n^-)/2$, remove every minority
term and every majority term whose endpoint has an outside
ancestor $t$ with $G(t)\le2B(t)$.  Denote the removed and kept
arrays by $\tau_n^{\mathrm r},\tau_n^{\mathrm k}$.
The first such stopping nodes have disjoint descendant sets.
The removed majority mass is at most $2\|b_n\|_1$; hence
$\|\tau_n^{\mathrm r}\|_1\le3d_n$.  This accounting is on
the labeled endpoint terms, before any cancellation in their
half-difference, so it bounds that difference as well.

At a node having a stopping ancestor the kept occupancy vanishes.
At every other outside node, $G(t)>2B(t)$ and
\[
 |(J_n\tau_n^{\mathrm k})(t)|
 \le G(t)\le2|G(t)-B(t)|=2|\bar h_n(t)|.
\]
Therefore
\[
 |Qy_n|_n\le\|Q\widetilde h_n\|_2
                  +2\|Q\bar h_n\|_2+6d_n.
\]
Lemma~\ref{qp:energy} and the triangle inequality in the outer
Euclidean sum give $\|Qy\|_{\mathrm q}\le9\sqrt E$.
Let $\eta\downarrow0$.
\end{proof}

\begin{corollary}\label{qp:lens-nine}
For every $\rho>0$, choose $0<\gamma<1/2$ such that
\[
 9\sqrt{(1+\gamma)^2-(1-2\gamma)^2}<\rho/3.
\]
If $\|x\|_{\mathrm q}\ge1-\gamma$, its symmetric lens
$\{y:\|x\pm y\|_{\mathrm q}\le1\}$ has no infinite
$\rho$-separated subset.
\end{corollary}

\begin{proof}
Approximate $x$ within $\gamma$ by a finite-support $x_0$, and
choose a finite ancestral head for it.  The tail of every point
in the lens has norm less than $\rho/3$ by
Proposition~\ref{qp:tail-nine}.  The points in the lens have norm
at most one, so in an infinite sequence two head projections are
less than $\rho/3$ apart.  Their full distance is less than $\rho$.
\end{proof}

The next formulation chooses one head from a nonzero center and a relative
endpoint excess before any displacement is supplied. It therefore controls
every admissible displacement at that center with the same head. This is
the form needed when several displacements are compared with one center.

\begin{proposition}\label{qp:relative-seventytwo}
For every nonzero $x\in X_{\mathrm q}$ and $0<\delta\le1$ there
is a finite ancestral head $A$, depending only on $x,\delta$, such
that for every $u\in X_{\mathrm q}$,
\[
 \|x+u\|_{\mathrm q},\ \|x-u\|_{\mathrm q}
       \le(1+\delta)\|x\|_{\mathrm q}
 \quad\Longrightarrow\quad
 \|Q_Au\|_{\mathrm q}\le72\sqrt\delta\,\|x\|_{\mathrm q}.
\]
\end{proposition}

\begin{proof}
Scale to $\|x\|_{\mathrm q}=1$.  Choose $x_0$ of finite node
support with $\|x-x_0\|_{\mathrm q}\le\delta$ and an ancestral
head $A$ containing its support.  Then
$\|x_0\pm u\|_{\mathrm q}\le1+2\delta$.  Choose representations
$x_0\pm u=h^\pm+J\lambda^\pm$ with each squared cost at most
$(1+3\delta)^2$.  Using the notation above for these representations,
let $C$ be their average squared cost and set
\[
\mathcal D=C-\left(\|P\bar h\|_H^2+\sum_nq_n^2\right).
\]
The parallelogram identity and the scalar square inequality used in
Lemma~\ref{qp:energy} give, before comparing the center representation
with its norm,
\[
 \begin{aligned}
 \|Q\bar h\|_H^2+\|\widetilde h\|_H^2+
       \sum_n(M_n^2-q_n^2)&\le\mathcal D,\\
 \tfrac12(\|Qh^+\|_H^2+\|Qh^-\|_H^2)+
       \sum_nd_n^2&\le\mathcal D.
 \end{aligned}
\]
The second inequality uses $M_n^2-q_n^2\ge d_n^2$ and the parallelogram
identity for $Qh^\pm$. The center representation has cost at least
$\|x_0\|_{\mathrm q}$, so
\begin{equation}\label{qp:seventytwo-energy}
 0\le\mathcal D\le(1+3\delta)^2-\|x_0\|_{\mathrm q}^2
 \le(1+3\delta)^2-(1-\delta)^2\le16\delta.
\end{equation}

Orient each nondiscarded fiber $C_r$ by the sign of $\nu_n(r)$, using
$+1$ for a zero folded coefficient or a discarded fiber.  This
time use the full, rather than
half, outside absolute masses: let $p_n^\pm$ agree with the
orientation and $b_n^\pm$ disagree with it, and put
$B_n=\|b_n^++b_n^-\|_1$.
The outside minority mass is at most the total cancellation
bound, giving $B_n\le2d_n$ and
\begin{equation}\label{qp:outside-bad}
 \sum_nB_n^2\le4\mathcal D.
\end{equation}
For outside nodes let
$G(t)=(p_n^++p_n^-)(S_t)$ and
$V(t)=(b_n^++b_n^-)(S_t)$.
Since the averaged represented vector vanishes there,
\[
 |G(t)-V(t)|=|h_n^+(t)+h_n^-(t)|.
\]
Stop at the first outside nodes satisfying $G(t)>0$ and
$V(t)\ge G(t)/2$.  Their descendant sets are disjoint, and
the total majority mass in them is at most $2B_n$.
Remove all these majority terms and all minority terms from the
outside part of $\lambda_n^+$; denote what remains by
$\lambda_n^{\mathrm k}$ and what was removed by
$\lambda_n^{\mathrm r}$.  Then
$\|\lambda_n^{\mathrm r}\|_1\le3B_n$ and, by
\eqref{qp:outside-bad},
\[
 \|\lambda^{\mathrm r}\|_L\le6\sqrt{\mathcal D},
 \qquad \|QJ\lambda^{\mathrm r}\|_{\mathrm q}\le12\sqrt{\mathcal D}.
\]
Every outside node either has a stopping ancestor, has $G(t)=0$,
or satisfies $V(t)<G(t)/2$.  In all cases the kept occupancy is
at most $2|G(t)-V(t)|$ in absolute value.  Thus
\[
 \|QJ\lambda^{\mathrm k}\|_H
 \le2\|Q(h^++h^-)\|_H\le4\sqrt{\mathcal D}.
\]
The atomic terms whose endpoints belong to $A$ have zero tail.
The representation of $x_0+u$, together with
$\|Qh^+\|_H\le\sqrt{2\mathcal D}$, now gives
\[
 \|Qu\|_{\mathrm q}\le(12+4+\sqrt2)\sqrt{\mathcal D}
 \le18\sqrt{\mathcal D}\le72\sqrt\delta.
\]
All stopping decisions use a first ancestor, so infinite
coefficient support and countably many children cause no change.
The choice of $A$ preceded the choice of $u$, as required.
\end{proof}

\subsection{From a component estimate to the Hilbert sum}
\label{qp:rec:section}

We now obtain a finite-codimensional midpoint estimate from normalized
component bounds. The proof applies a local lens estimate on components
with controlled relative endpoint excess and estimates the remaining
components by their contribution to the total squared excess.

\begin{lemma}[Component estimate]\label{qp:rec:component}
Let $A\subset T_n$ be a finite nonempty ancestral head, with
$P_Ax=x$, $P_Ay=0$, and $|x|_n=1$.  If $0\le e\le1$ and
$|x\pm y|_n\le1+e$, then
\begin{equation}\label{qp:rec:component-bound}
 |y|_n\le(6+\sqrt2)\sqrt{(1+e)^2-1}\le13\sqrt e.
\end{equation}
\end{lemma}

\begin{proof}
Apply the proof of Theorem~\ref{thm:quadratic-tail} in this single
component, retaining the coefficient $6+\sqrt2$ before its final
rounding to $8$.  Here $Q_Ay=y$.  Explicitly, approximate endpoint
representations of cost below $1+e+\varepsilon$ give
$\kappa=(1+e+\varepsilon)^2-1$, Hilbert tails at most
$\sqrt{2\kappa}$, combined Hilbert tail at most $2\sqrt\kappa$,
and folded-mass loss at most $\sqrt\kappa$ by
Lemma~\ref{qp:energy}.  Because the head contains the root, no
fiber is discarded: the combined minority mass is exactly this
folded-mass loss.  Grouping by the common truncated endpoint and
clipping one majority occupancy against the combined minority
occupancy gives two signed stopped flows of norm at most
$2\sqrt\kappa$ each, by Lemma~\ref{qp:flow}.  The remaining
majority occupancy is bounded in Hilbert norm by
$2\sqrt\kappa$.  Adding the original Hilbert tail gives
$(6+\sqrt2)\sqrt\kappa$.  These are the same decomposition and
signed path representations used in that theorem, so no invariance
of $|\cdot|_n$ under arbitrary coordinate sign changes is needed.
Let $\varepsilon\downarrow0$.  Finally
$(1+e)^2-1\le3e$ and
$3(6+\sqrt2)^2=114+36\sqrt2<169$.
\end{proof}

The component estimate controls a lens after normalization by the center
norm.  In a Hilbert sum that norm can be very small or zero.  The next
argument separates the components carrying a large relative endpoint
excess before making any such division.

\begin{proposition}[Finite centers in the Hilbert sum]
\label{qp:rec:finite-center}
Let \(x\in X_{\mathrm q}\) be a unit vector with finite coordinate
support.  For each \(n\) with \(x_n\ne0\), choose a finite nonempty
ancestor-closed \(A_n\subset T_n\) containing its support, and set
\[
 Y_x=\{y\in X_{\mathrm q}:P_{A_n}y_n=0
                           \text{ whenever }x_n\ne0\}.
\]
Then \(Y_x\) is a closed finite-codimensional subspace.  If
\(0<e\le1/8\), \(y\in Y_x\), and
\(\|x+y\|_{\mathrm q},\|x-y\|_{\mathrm q}\le1+e\), then
\begin{equation}\label{qp:rec:outer-bound}
 \|y\|_{\mathrm q}^2\le\frac{519}{2}e^{1/3},
 \qquad \|y\|_{\mathrm q}\le17e^{1/6}.
\end{equation}
\end{proposition}

\begin{proof}
The conditions defining \(Y_x\) involve only finitely many continuous
coordinate functionals, so the assertion about codimension follows.
Put
\[
 a_n=|x_n|_n,\quad b_n=|x_n+y_n|_n,\quad c_n=|x_n-y_n|_n,
 \quad v_n=\frac{b_n^2+c_n^2}{2},\quad d_n=v_n-a_n^2.
\]
The triangle inequality gives
\(a_n,|y_n|_n\le(b_n+c_n)/2\).  Therefore
\begin{equation}\label{qp:rec:outer-energy}
 d_n\ge0,\qquad |y_n|_n^2\le v_n,
 \qquad \sum_na_n^2=1,
 \qquad \sum_nd_n\le(1+e)^2-1\le3e.
\end{equation}
Fix \(0<\eta\le1\).  Write
\[
 E_\eta=\{n:d_n>\eta a_n^2\}.
\]
On these components \(v_n=a_n^2+d_n\le(1+1/\eta)d_n\), and hence
\begin{equation}\label{qp:rec:exceptional}
 \sum_{n\in E_\eta}|y_n|_n^2
 \le3e(1+1/\eta).
\end{equation}
If \(a_n=0\) and \(y_n\ne0\), then
\(d_n=v_n=|y_n|_n^2>0\), so \(n\in E_\eta\).
If both vectors vanish in that component, its contribution is zero.
Thus every remaining nonzero contribution has \(a_n>0\).

Now assume additionally that \(\eta\le1/4\).  For
\(n\notin E_\eta\) with \(a_n>0\), let
\(m_n=(b_n+c_n)/2\).  We have \(m_n\ge a_n\) and
\[
 m_n^2+\left(\frac{b_n-c_n}{2}\right)^2=v_n
                                             \le(1+\eta)a_n^2.
\]
It follows that
\[
 \frac{\max(b_n,c_n)}{a_n}
 \le\sqrt{1+\eta}+\sqrt\eta=1+\delta_\eta,
 \qquad
 \delta_\eta:=\sqrt{1+\eta}+\sqrt\eta-1
                  \le\frac32\sqrt\eta\le1.
\]
Apply Lemma~\ref{qp:rec:component} to
\(x_n/a_n\) and \(y_n/a_n\).  Summing its squared estimate gives
\begin{equation}\label{qp:rec:regular}
 \sum_{n\notin E_\eta}|y_n|_n^2
 \le169\delta_\eta\sum_{n\notin E_\eta}a_n^2
 \le\frac{507}{2}\sqrt\eta.
\end{equation}
Take \(\eta=e^{2/3}\), which is at most \(1/4\) when \(e\le1/8\).
Combining \eqref{qp:rec:exceptional} and \eqref{qp:rec:regular},
\[
 \|y\|_{\mathrm q}^2
 \le\frac{507}{2}e^{1/3}+3e+3e^{1/3}
 \le\frac{519}{2}e^{1/3}.
\]
The second assertion follows from \(519/2<17^2\).
\end{proof}

\begin{corollary}[Completed centers]\label{qp:rec:all-centers}
For every \(x\in X_{\mathrm q}\) with \(\|x\|_{\mathrm q}=1\) and
\(0<e\le1/16\), there is a closed finite-codimensional subspace
\(Y\subset X_{\mathrm q}\) such that
\begin{equation}\label{qp:rec:complete-bound}
 y\in Y,\quad
 \|x+y\|_{\mathrm q},\|x-y\|_{\mathrm q}\le1+e
 \quad\Longrightarrow\quad
 \|y\|_{\mathrm q}\le17(2e)^{1/6}.
\end{equation}
In particular, for every \(t>0\) there is an \(e_t>0\), independent
of the unit center \(x\), for which some closed finite-codimensional
\(Y\) satisfies
\[
 \inf_{u\in S_Y}
 \max\{\|x+tu\|_{\mathrm q},\|x-tu\|_{\mathrm q}\}\ge1+e_t.
\]
Here \(S_Y=\{u\in Y:\|u\|_{\mathrm q}=1\}\).  Thus the maximum
asymptotic midpoint modulus is positive at every positive radius.
\end{corollary}

\begin{proof}
By finite-coordinate density, choose a finitely supported unit vector
\(x^0\) with \(\|x-x^0\|_{\mathrm q}<e\).  Such unit vectors are
dense: approximate \(x\) by a nonzero finitely supported vector and
normalize it.  Use the finite coordinate exclusions for \(x^0\) from
Proposition~\ref{qp:rec:finite-center}.  If \(y\) satisfies the
hypotheses of \eqref{qp:rec:complete-bound}, then
\[
 \|x^0+y\|_{\mathrm q},\|x^0-y\|_{\mathrm q}\le1+2e.
\]
Since \(2e\le1/8\), Proposition~\ref{qp:rec:finite-center} proves
\eqref{qp:rec:complete-bound}.  Given \(t>0\), choose
\(0<e_t\le1/16\) with \(17(2e_t)^{1/6}<t\).
For \(u\in S_Y\), both endpoint norms at \(tu\) cannot be at most
\(1+e_t\), by \eqref{qp:rec:complete-bound}.  Each corresponding
maximum is therefore greater than \(1+e_t\), giving the displayed
infimum bound.
\end{proof}
